\begin{lemma}
Let $\mathcal H$ be a $3$-uniform hypergraph of maximum degree at most $D$,
let $e\in E(\mathcal H)$ intersect exactly $3(D-1)$ other edges, and let $G$
be the subgraph of $L(\mathcal H)$ induced on $N_{L(\mathcal H)}(e)$.  With
$X=X(e,\mathcal H)$ and
\[
w_{\mathcal H,e}(X)=
\sum_{x\in X}\sum_{v<v'\in e}
\deg_{\mathcal H}(v,x)\deg_{\mathcal H}(v',x),
\]
the number $P(e)$ of independent pairs in $N_{L(\mathcal H)}(e)$ satisfies
\[
P(e)\ge 3D^2-6D+3-w_{\mathcal H,e}(X).
\]
\end{lemma}
\begin{proof}
By \noderef{LineGraphOfHypergraph}, the vertices of the line graph are the
edge indices of $\mathcal H$, and the neighborhood $N_{L(\mathcal H)}(e)$ is
exactly the set of edge indices $h\ne e$ whose hyperedges meet $e$.  The
hypothesis says that this set has cardinality $3(D-1)=3D-3$.  Therefore the
total number of unordered two-element subsets of the neighborhood is
\[
\binom{3D-3}{2}.
\]
By \noderef{IndependentPairCount}, $P(e)$ counts exactly the two-element
subsets of $N_{L(\mathcal H)}(e)$ whose two vertices are not adjacent in the
line graph.  Since $G$ is the subgraph induced on the same neighborhood, the
edges of $G$ are exactly the adjacent two-element subsets of this
neighborhood.  The adjacent and non-adjacent two-element subsets partition
all two-element subsets, so
\[
P(e)=\binom{3D-3}{2}-|E(G)|.
\]

The edge-count identity \noderef{SaturatedLineGraphEdgeCount}, applied to
this $3$-uniform hypergraph of maximum degree at most $D$ and to the edge
$e$ that meets exactly $3(D-1)$ other edges, gives
\[
|E(G)|=3\binom{D-1}{2}+w_{\mathcal H,e}(X)-Y_{\mathcal H}(e),
\]
and the same cited statement gives $Y_{\mathcal H}(e)\ge0$.  Substituting
this identity into the previous display gives
\[
P(e)=
\binom{3D-3}{2}
-3\binom{D-1}{2}
-w_{\mathcal H,e}(X)+Y_{\mathcal H}(e).
\]
Since $Y_{\mathcal H}(e)\ge0$, discarding it can only decrease the
right-hand side, and hence
\[
P(e)\ge
\binom{3D-3}{2}-3\binom{D-1}{2}-w_{\mathcal H,e}(X).
\]
The elementary binomial calculation
\[
\binom{3D-3}{2}-3\binom{D-1}{2}=3D^2-6D+3
\]
then gives the asserted lower bound.
\end{proof}
